\chapter*{Заключение}                       % Заголовок
\addcontentsline{toc}{chapter}{Заключение}  % Добавляем его в оглавление

%% Согласно ГОСТ Р 7.0.11-2011:
%% 5.3.3 В заключении диссертации излагают итоги выполненного исследования, рекомендации, перспективы дальнейшей разработки темы.
%% 9.2.3 В заключении автореферата диссертации излагают итоги данного исследования, рекомендации и перспективы дальнейшей разработки темы.
%% Поэтому имеет смысл сделать эту часть общей и загрузить из одного файла в автореферат и в диссертацию:

Основные результаты работы заключаются в следующем.
%% Согласно ГОСТ Р 7.0.11-2011:
%% 5.3.3 В заключении диссертации излагают итоги выполненного исследования, рекомендации, перспективы дальнейшей разработки темы.
%% 9.2.3 В заключении автореферата диссертации излагают итоги данного исследования, рекомендации и перспективы дальнейшей разработки темы.
\begin{enumerate}
  \item При создании гибридного источника излучения на длине волны 3 мкм с волоконной накачкой в маломодовом ОВ доставки впервые обнаружено взаимное влияние ОМЧВС и ММЧВС, имеющих общую антистоксовую компоненту на длине волны 1003~нм. Показано, что синхронизация импульсов накачки (1030~нм) и затравки (1562~нм) увеличивает спектральную мощность компонент ММЧВС на 20 дБ, что впоследствии снижает эффективность ГРЧ в н-о кристалле. Предложен и реализован метод подавления взаимного влияния ЧВС с помощью модового фильтра на входе ОВ, позволяющий увеличить допустимую длину ОВ доставки с 2.5 м до 3.0 м без существенной потери полезной мощности накачки.
  \item Впервые экспериментально исследована эволюция КФС при фотоиндуцированной ГВГ в германосиликатном ОВ в режиме самозатравки. Показано, что в процессе оптического полинга ширина КФС и их сдвиг монотонно уменьшаются до стационарных значений, а форма кривых сохраняет выраженную асимметрию. Проведён соответствующий численный расчёт с использованием модели~\cite{antonyuk2001self} с добавлением учёта керровского эффекта. Обнаружено, что на поздних стадиях эффективность ГВГ остаётся практически постоянной вследствие динамического равновесия между уменьшением спектральной расстройки и ростом пиковой эффективности преобразования. Полученные результаты создают основу для разработки физической модели фотоиндуцированной ГВГ в ОВ, которая позволит прогнозировать уровень паразитной ГВГ в ОВ доставки мощных лазеров, в которых данный эффект может быть нежелательным, поскольку он может приводить к снижению порога оптического разрушения н-о кристаллов при ГРЧ в гибридных лазерных системах.
  \item Предложены поправочные выражения для трёх методик обработки температурных кинетик при разогреве образца лазерным излучением в методе ЛК, позволяющие компенсировать нестабильность температуры окружающей среды и коэффициента теплоотдачи. Экспериментально показано, что применение поправок снижает разброс измеренных значений коэффициента оптического поглощения более чем в два раза (с 23\% до 11\%).
\end{enumerate}

%И какая-нибудь заключающая фраза.

%Последний параграф может включать благодарности.

\subsubsection*{Благодарности}
Кого благодарим:
Коняшкина Алексея Викторовича - за то что связал три несвязных между собой темы и поддерживал меня что эту херь можно защитить и что она норм написана ну типо вселил веру в меня но не так пафосно

Терещенко Никиту Вячеславовича, Зотова Кирилла Вадимовича - за высокий интерес к науке, многолетнюю плодотворную совместную работу с офигенно эффективным разделением обязанностей, с обсуждением науки как на семинаре так и за бутылкой пива, за фактически сложившийся научный мини-коллектив

Цыпкина Виктора Павловича профессионализм при проведении экспериментов, постоянную готовность помочь, за деловой подход и за те же обсуждения рабочих и научных вопросов с пивом, что сохраняло мотивацию не останавливаться.

Тыртышного Валентина Александровича, Баранова Андрея Игоревича, Ларионова Игоря Александровича за мой непрерывный профессиональный и научный рост, без которого диссертация была бы невозможна. 

Суровцеву Викторию Павловну, Казаринову Дарью Дмитриевну за неоценимую (синоним без пафоса!!) помощь при проведении экспериментов

Коллектив кафедры фотоники: Мясникова Даниила Владимировича, Шайдуллина Рената Ильгизовича, Коняшкина Алексея Викторовича (или в второй раз не надо?), 
%Автор выражает глубокую признательность своему научному руководителю, кандидату физико-математических наук **Коняшкину Алексею Викторовичу**, за постоянное внимание к работе, плодотворные обсуждения материалов диссертации, ценные рекомендации при анализе экспериментальных и теоретических результатов, а также за создание атмосферы, способствующей глубокому научному поиску.
%
%Автор искренне благодарит **Рябушкина Олега Алексеевича** и весь преподавательский состав кафедры фотоники МФТИ (НИУ) за важные замечания при подготовке работы к защите и формирование профессиональных компетенций, необходимых для выполнения данной диссертации.
%
Автор приносит благодарность компании ООО «ВПГ Лазеруан» за предоставленную возможность проведения экспериментальных исследований с использованием современной приборной и компонентной базы.
%
%Автор высоко ценит содействие и помощь коллег: **Зотова К.В., Алояна Г., Шебаршиной И.В., Пакулевой С.В., Терещенко Н.В., Грищенко И.В., Вершинина О.И., Никитина Д.Г., Тыртышного В.А., Гуляшко А.А., Ларионова И.С.** — за сотрудничество при выполнении экспериментальных исследований, обсуждение полученных результатов и создание творческой рабочей атмосферы. Отдельная благодарность **Андрееву А.О., Андреевой В.А., Мясникову Д.В., Евтихиеву Н.Н.** за поддержку и ценные советы на различных этапах работы.



Наконец, автор выражает сердечную благодарность своим близким и родным за неизменную моральную поддержку, терпение и веру в успех на протяжении всего периода подготовки диссертации.

 
%В заключение автор выражает благодарность и большую признательность научному руководителю
%Иванову~И.\,И. за поддержку, помощь, обсуждение результатов и~научное
%руководство. Также автор благодарит Сидорова~А.\,А. и~Петрова~Б.\,Б.
%за помощь в~работе с~образцами, Рабиновича~В.\,В. за предоставленные
%образцы и~обсуждение результатов, Занудятину~Г.\,Г. и авторов шаблона
%*Russian-Phd-LaTeX-Dissertation-Template* за~помощь в оформлении
%диссертации. Автор также благодарит много разных людей
%и~всех, кто сделал настоящую работу автора возможной.
